% GENERATED FILE. Edit canonical lessons, then run npm run generate:books.
% canonical-source-hash: fnv1a64:3e3cfbe3396e103f
% canonical-lessons: LA-C239-simul, LA-C239-fortasse, LA-C239-quoque, LA-C239-quasi, LA-C239-propter

\chapter{At the same time, Perhaps, Also, As if, Because of}
\label{ch:la-at-the-same-time-perhaps-also-as-if-because-of}
\hlchaptermodality{239}

\begin{chapteropening}
\textbf{By the end of this chapter:} \emph{I can say at the same time, perhaps, also, as if and because of.}

Complete the last lesson of chapter 239: I can say at the same time, perhaps, also, as if and because of.
\end{chapteropening}

\section[simul]{simul — at the same time}

\label{lesson:LA-C239-simul}



\begin{quote}
Before the new one: say the Latin for face to face, then the Latin for by which way.
\end{quote}

\subsection*{You'll want to know: simul}
\textbf{simul} — \textquotedblleft{}at the same time\textquotedblright{}.

\begin{cousinweb}
\textbf{simul} is inherited, usually traced to the Proto-Indo-European root \emph{*sem-}, \textquotedblleft{}one, together\textquotedblright{}. Greek \emph{hama}, \textquotedblleft{}together\textquotedblright{}, is a cousin (Lewis \& Short), and so is English \emph{same}.
\end{cousinweb}

\begin{grammarlens}[title={What you've built}]
One more word for this chapter.
\end{grammarlens}

\subsection*{Your turn}
\begin{itemize}
\raggedright
  \item \emph{Say it:} \emph{simul}
  \item \emph{Say it:} \emph{simul}, once more
  \item \emph{Say it:} say \emph{simul}
  \item \emph{Recall:} say the Latin for face to face, then the Latin for by which way, then say \emph{simul} again
\end{itemize}

\begin{tcolorbox}[breakable,colback=teal!4,colframe=teal!35!black,title={Before you move on}]
What does simul mean? (\textquotedblleft{}At the same time\textquotedblright{}.) Say it once more.
\end{tcolorbox}
\section[fortasse]{fortasse — perhaps}

\label{lesson:LA-C239-fortasse}



\begin{quote}
Before the new one: say the Latin for by which way, then the Latin for at the same time.
\end{quote}

\subsection*{You'll want to know: fortasse}
\textbf{fortasse} — \textquotedblleft{}perhaps\textquotedblright{}.

\begin{cousinweb}
\textbf{fortasse} is usually explained as built on \emph{forte}, \textquotedblleft{}by chance\textquotedblright{}, and \emph{an}, \textquotedblleft{}whether\textquotedblright{} (Lewis \& Short): perhaps.
\end{cousinweb}

\begin{grammarlens}[title={What you've built}]
One more word for this chapter.
\end{grammarlens}

\subsection*{Your turn}
\begin{itemize}
\raggedright
  \item \emph{Say it:} \emph{fortasse}
  \item \emph{Say it:} \emph{fortasse}, once more
  \item \emph{Say it:} say \emph{fortasse}
  \item \emph{Recall:} say the Latin for by which way, then the Latin for at the same time, then say \emph{fortasse} again
\end{itemize}

\begin{tcolorbox}[breakable,colback=teal!4,colframe=teal!35!black,title={Before you move on}]
What does fortasse mean? (\textquotedblleft{}Perhaps\textquotedblright{}.) Say it once more.
\end{tcolorbox}
\section[quoque]{quoque — also, too}

\label{lesson:LA-C239-quoque}



\begin{quote}
Before the new one: say the Latin for at the same time, then the Latin for perhaps.
\end{quote}

\subsection*{You'll want to know: quoque}
\textbf{quoque} — \textquotedblleft{}also, too\textquotedblright{}.

\begin{cousinweb}
\textbf{quoque} has no sure origin; Lewis \& Short call it doubtful. It is a different word from \emph{quōque}, \textquotedblleft{}and by which\textquotedblright{}.
\end{cousinweb}

\begin{grammarlens}[title={What you've built}]
One more word for this chapter.
\end{grammarlens}

\subsection*{Your turn}
\begin{itemize}
\raggedright
  \item \emph{Say it:} \emph{quoque}
  \item \emph{Say it:} \emph{quoque}, once more
  \item \emph{Say it:} say \emph{quoque}
  \item \emph{Recall:} say the Latin for at the same time, then the Latin for perhaps, then say \emph{quoque} again
\end{itemize}

\begin{tcolorbox}[breakable,colback=teal!4,colframe=teal!35!black,title={Before you move on}]
What does quoque mean? (\textquotedblleft{}Also, too\textquotedblright{}.) Say it once more.
\end{tcolorbox}
\section[quasi]{quasi — as if}

\label{lesson:LA-C239-quasi}



\begin{quote}
Before the new one: say the Latin for perhaps, then the Latin for also.
\end{quote}

\subsection*{You'll want to know: quasi}
\textbf{quasi} — \textquotedblleft{}as if\textquotedblright{}.

\begin{cousinweb}
\textbf{quasi} is contracted from \emph{quam sī}, \textquotedblleft{}as if\textquotedblright{} (Lewis \& Short).
\end{cousinweb}

\begin{grammarlens}[title={What you've built}]
One more word for this chapter.
\end{grammarlens}

\subsection*{Your turn}
\begin{itemize}
\raggedright
  \item \emph{Say it:} \emph{quasi}
  \item \emph{Say it:} \emph{quasi}, once more
  \item \emph{Say it:} say \emph{quasi}
  \item \emph{Recall:} say the Latin for perhaps, then the Latin for also, then say \emph{quasi} again
\end{itemize}

\begin{tcolorbox}[breakable,colback=teal!4,colframe=teal!35!black,title={Before you move on}]
What does quasi mean? (\textquotedblleft{}As if\textquotedblright{}.) Say it once more.
\end{tcolorbox}
\section[propter]{propter — because of}

\label{lesson:LA-C239-propter}



\begin{quote}
Before the new one: say the Latin for also, then the Latin for as if.
\end{quote}

\subsection*{You'll want to know: propter}
\textbf{propter} — \textquotedblleft{}because of\textquotedblright{}.

\begin{cousinweb}
\textbf{propter} is contracted from an older \emph{*propiter}, built on \emph{prope}, \textquotedblleft{}near\textquotedblright{} (Lewis \& Short): near by, then because of.
\end{cousinweb}

\begin{grammarlens}[title={What you've built}]
One more word for this chapter.
\end{grammarlens}

\subsection*{Your turn}
\begin{itemize}
\raggedright
  \item \emph{Say it:} \emph{propter}
  \item \emph{Say it:} \emph{propter}, once more
  \item \emph{Say it:} say \emph{propter}
  \item \emph{Recall:} say the Latin for also, then the Latin for as if, then say \emph{propter} again
\end{itemize}

\begin{tcolorbox}[breakable,colback=teal!4,colframe=teal!35!black,title={Before you move on}]
What does propter mean? (\textquotedblleft{}Because of\textquotedblright{}.) Say it once more.
\end{tcolorbox}
